\documentclass[11pt]{article}
\usepackage[T1]{fontenc}
\usepackage[utf8]{inputenc}
\usepackage{newpxtext}
\usepackage[a4paper,margin=27mm,headheight=15pt]{geometry}
\usepackage{microtype}
\usepackage{amsmath,amsthm,mathtools}
\usepackage{newpxmath}
\usepackage{booktabs,array,longtable}
\usepackage{xcolor}
\usepackage{enumitem}
\usepackage{fancyhdr}
\usepackage{listings}
\usepackage{hyperref}
\usepackage{aliascnt}
\usepackage[nameinlink,noabbrev]{cleveref}
\definecolor{Ink}{HTML}{20354C}
\definecolor{Accent}{HTML}{087F8C}
\definecolor{Pale}{HTML}{F2F6F8}
\hypersetup{colorlinks=true,linkcolor=Ink,citecolor=Accent,urlcolor=Accent,
  pdftitle={Guarded periodic blocks and maximal order types below omega squared},
  pdfauthor={Research report generated with ChatGPT},
  pdfsubject={A proposed solution to Altman's finite-alphabet sharpness question}}
\pagestyle{fancy}
\fancyhf{}
\fancyhead[L]{\small\scshape Guarded periodic blocks}
\fancyhead[R]{\small\scshape Words below $\omega^2$}
\fancyfoot[C]{\thepage}
\renewcommand{\headrulewidth}{0.3pt}
\setlength{\parskip}{3pt}
\setlength{\parindent}{1.2em}
\setlist{itemsep=3pt,topsep=5pt}
\emergencystretch=2em
\numberwithin{equation}{section}
\newtheorem{theorem}{Theorem}[section]
\theoremstyle{plain}
\newaliascnt{lemma}{theorem}
\newtheorem{lemma}[lemma]{Lemma}
\aliascntresetthe{lemma}
\crefname{lemma}{lemma}{lemmas}
\Crefname{lemma}{Lemma}{Lemmas}
\newaliascnt{proposition}{theorem}
\newtheorem{proposition}[proposition]{Proposition}
\aliascntresetthe{proposition}
\crefname{proposition}{proposition}{propositions}
\Crefname{proposition}{Proposition}{Propositions}
\newaliascnt{corollary}{theorem}
\newtheorem{corollary}[corollary]{Corollary}
\aliascntresetthe{corollary}
\crefname{corollary}{corollary}{corollaries}
\Crefname{corollary}{Corollary}{Corollaries}
\theoremstyle{definition}
\newaliascnt{definition}{theorem}
\newtheorem{definition}[definition]{Definition}
\aliascntresetthe{definition}
\crefname{definition}{definition}{definitions}
\Crefname{definition}{Definition}{Definitions}
\newaliascnt{example}{theorem}
\newtheorem{example}[example]{Example}
\aliascntresetthe{example}
\crefname{example}{example}{examples}
\Crefname{example}{Example}{Examples}
\theoremstyle{remark}
\newaliascnt{remark}{theorem}
\newtheorem{remark}[remark]{Remark}
\aliascntresetthe{remark}
\crefname{remark}{remark}{remarks}
\Crefname{remark}{Remark}{Remarks}

\newcommand{\ot}{\operatorname{o}}
\newcommand{\J}{\mathcal J}
\newcommand{\F}{\mathcal F}
\newcommand{\W}{\mathcal W}
\newcommand{\E}{\mathcal E}
\newcommand{\HH}{H}
\newcommand{\N}{\mathbb N}
\newcommand{\emp}{\varepsilon}
\newcommand{\len}{\operatorname{len}}
\newcommand{\Max}{\operatorname{Max}}
\newcommand{\rec}{\operatorname{Rec}}
\newcommand{\NF}{\operatorname{NF}}
\newcommand{\down}{\mathord\downarrow}
\newcommand{\sub}{\preccurlyeq}
\newcommand{\nsub}{\npreccurlyeq}
\newcommand{\equ}{\sim}
\newcommand{\lex}{\mathrm{lex}}
\lstset{basicstyle=\ttfamily\small,columns=fullflexible,keepspaces=true,
  frame=single,rulecolor=\color{Ink!25},backgroundcolor=\color{Pale},
  breaklines=true,showstringspaces=false,aboveskip=9pt,belowskip=9pt}

\title{\vspace{-9mm}{\color{Ink}\bfseries Guarded periodic blocks and maximal order types\par\medskip
of words below $\boldsymbol{\omega^2}$}\\[9pt]
\large A proposed resolution of Altman's two-letter sharpness question}
\author{Research report generated with ChatGPT}
\date{19 September 2026}

\begin{document}
\maketitle
\begin{abstract}
Let $X$ be a finite poset, and order the $X$-labelled words of ordinal length less than $\omega^2$ by strictly increasing, label-monotone subsequence embeddings. We give a proof that, for $n=|X|\geq2$,
\[
 \boxed{\displaystyle
 \ot\bigl(s_{\omega^2}(X)\bigr)
   =\omega^{\omega^{\,n+|\J(X)|-2}} ,}
\]
where $\J(X)$ contains all downward-closed subsets, including the empty and full subsets. In particular, the two-letter antichain has maximal order type $\omega^{\omega^4}$, attaining the upper bound whose sharpness was left unverified in Example~3.15 of Altman's March 2026 revision. The two-letter chain has type $\omega^{\omega^3}$.

The main step constructs order embeddings of all finite Cartesian powers by means of periodic separators. A terminal letter protects components from a proper recurrent ideal; an infinite terminal guard handles full support. We also classify arbitrary allowed families of recurrent ideals, identify an essential exceptional family, compute exact-length strata, and give an exact symbolic embedding algorithm. The accompanying programs perform exhaustive binary checks and additional tests on small finite posets. These finite tests audit the constructions; the transfinite conclusions depend on the proofs.
\end{abstract}

\begin{center}
\fcolorbox{Ink!25}{Pale}{\parbox{0.91\linewidth}{\small
\textbf{Status and scope.} This is a proof manuscript, not a peer-reviewed or proof-assistant-verified result. The explicit unresolved question was checked in the cited source. The claim of a solution rests on the argument given here, not on a claim of exhaustive literature coverage. The theorem concerns \emph{strictly increasing} embeddings and the strict length cutoff $\omega^2$.}}
\end{center}
\vfill
\noindent\textbf{Keywords.} Well-quasi-orders; maximal order type; transfinite words; recurrent ideals; ordinal arithmetic; Higman order.
\newpage
\begingroup
\small
\setlength{\parskip}{0pt}
\tableofcontents
\endgroup
\newpage

\section{The problem and the result}\label{sec:intro}

A finite alphabet gives rise to two quite different sorts of complexity. There are arbitrarily long finite words, and there are infinite blocks in which specified letters recur without end. Below $\omega^2$, only finitely many such infinite blocks occur in any individual word. Nevertheless, allowing unboundedly many blocks produces a substantial increase in maximal order type.

The starting point is a concrete question in Altman~\cite[Example~3.15]{Altman}. For a two-element antichain $A_2$, the established upper bound is
\begin{equation}\label{eq:altman}
 \ot\bigl(s^F_{\omega^2}(A_2)\bigr)\leq\omega^{\omega^4}.
\end{equation}
The March 2026 revision leaves its sharpness unverified. The same example gives the improved upper bound $\omega^{\omega^3}$ for the two-element chain. The superscript $F$ means finite image; for a finite alphabet it is automatic, and we omit it.

We attack the missing lower bound. The construction turns out to work for every finite poset, and not just for antichains.

\begin{theorem}[Finite-alphabet formula]\label{thm:main}
Let $X$ be a finite poset. Put $n=|X|$ and $j=|\J(X)|$, where $\J(X)$ is the family of all downward-closed subsets of $X$. Then
\begin{equation}\label{eq:main}
 \ot\bigl(s_{\omega^2}(X)\bigr)=
 \begin{cases}
  1,&n=0,\\
  \omega^2,&n=1,\\
  \omega^{\omega^{n+j-2}},&n\geq2.
 \end{cases}
\end{equation}
\end{theorem}

The two small exceptional alphabet sizes are part of the theorem, not removable technicalities. In particular, substituting $n=1$ into the third line would give an incorrect answer.

For an antichain, every subset is downward closed. For a chain, downward-closed subsets are precisely the initial segments. Consequently the theorem specializes to
\begin{align}
 \ot\bigl(s_{\omega^2}(A_n)\bigr)
   &=\omega^{\omega^{n+2^n-2}} &&(n\geq2),\label{eq:anti}\\
 \ot\bigl(s_{\omega^2}(C_n)\bigr)
   &=\omega^{\omega^{2n-1}} &&(n\geq2).\label{eq:chain}
\end{align}
Thus \eqref{eq:altman} is attained, and the improved chain bound is attained as well.

\paragraph{What is being claimed as the contribution?}
The normal-form encoding and the finite Higman maximal-order-type formula are established ingredients. The central contribution of this manuscript is the guarded product construction in \cref{sec:guards,sec:product}, which makes the generator-count upper bound sharp. The support-family classification and the finite-poset formula follow from that construction. The exact-length calculations are supplementary consequences, not the foundation of the main lower bound.

\paragraph{Related work and the scope of the literature check.}
The classical maximal-linear-extension theory is due to de Jongh and Parikh~\cite{DJP}; the finite-word calculation is part of the de Jongh--Parikh--Schmidt theory~\cite{DJP,Schmidt}. Related recent work of Chopra and Pakhomov~\cite{CP} studies the limit cutoff $\omega^\omega$ through finite trees. Here we address the finite cutoff $\omega^2$ and the explicitly identified two-letter example. No claim is made that the present literature check excludes every earlier formulation of the general finite-poset formula.

\section{Conventions and classical ingredients}\label{sec:prelim}

\subsection{Words and the precise embedding relation}

\begin{definition}
An $X$-word $u$ of length $\alpha$ is a function $u:\alpha\to X$. We write $\len(u)=\alpha$. For words $u:\alpha\to X$ and $v:\beta\to X$, put $u\sub v$ when there is a strictly increasing function
\[
 f:\alpha\longrightarrow\beta
 \quad\text{such that}\quad
 u(\xi)\leq_X v(f(\xi))\qquad(\xi<\alpha).
\]
There is no continuity requirement on $f$. Write $u\equ v$ for mutual embeddability. The empty word is $\emp$.
\end{definition}

The quasiorder $s_{\omega^2}(X)$ consists of all these words with length strictly less than $\omega^2$. Concatenation is ordinal concatenation. In particular,
\[
 k+\omega=\omega,
 \qquad
 \omega+\omega=\omega\cdot2,
 \qquad
 \omega r+k<\omega^2\quad(r,k\in\N).
\]
Finite material preceding an infinite block may therefore disappear from the \emph{length expression}, although its labels can still matter for embeddability.

Strictness of the index map is indispensable. Under a merely nondecreasing index map, an infinite constant source could be collapsed, and the results below would describe a different problem.

\subsection{Maximal order type}

A quasiorder $P$ is a well-quasi-order (wqo) when every infinite sequence $x_0,x_1,\ldots$ has $i<j$ with $x_i\leq_P x_j$. Its antisymmetric quotient is denoted $P/{\equ}$. All linear extensions and order embeddings below are understood on this quotient when necessary.

The maximal order type $\ot(P)$ is the largest order type of a linear extension of $P/{\equ}$. The existence of a largest such order type, rather than merely a supremum, is the maximal-linear-extension theorem~\cite{DJP}.

We use the following standard consequences. If $P$ is an induced suborder of $Q$, or more generally order-embeds into $Q$, then
\begin{equation}\label{eq:mono}
 \ot(P)\leq\ot(Q).
\end{equation}
A monotone surjection $Q\to P$ also implies \eqref{eq:mono}. For the latter assertion, choose a linear extension of $P$ of type $\ot(P)$ and order the fibres of the surjection consecutively according to this extension. Linearize each fibre. This gives a linear extension of $Q$ with type at least $\ot(P)$. The statement applies equally to a surjection onto equivalence classes.

The standard Cartesian-product and disjoint-union identities are
\begin{equation}\label{eq:naturalops}
 \ot(P\times Q)=\ot(P)\otimes\ot(Q),
 \qquad
 \ot(P\amalg Q)=\ot(P)\oplus\ot(Q),
\end{equation}
where $\oplus$ and $\otimes$ are the natural, or Hessenberg, operations~\cite{DJP}. These identities are used only for the supplementary exact-length calculations. The main theorem needs a weaker and elementary product lower bound, proved in \cref{lem:power}.

\subsection{The finite Higman formula}

For a finite nonempty poset of cardinality $m$, the order of finite words has maximal order type
\begin{equation}\label{eq:higman}
 \ot(X^*)=\HH_m,
 \qquad \HH_m:=\omega^{\omega^{m-1}}.
\end{equation}
This classical formula does not depend on the shape of the finite poset~\cite{DJP,Schmidt}. Its wqo assertion is Higman's lemma~\cite{Higman}. We invoke \eqref{eq:higman} as an established theorem, not as a result of the symbolic tests.

For orientation,
\[
 \HH_1=\omega,
 \quad \HH_2=\omega^\omega,
 \quad \HH_3=\omega^{\omega^2},
 \quad \HH_4=\omega^{\omega^3}.
\]
Ordinary ordinal operations are used unless $\oplus$ or $\otimes$ is explicitly written. The elementary identity driving the induction is
\begin{equation}\label{eq:amplification}
 \sup_{r\geq1}\HH_q^{\,r}
 =\sup_{r\geq1}\omega^{\omega^{q-1}\cdot r}
 =\omega^{\omega^q}
 =\HH_{q+1}\qquad(q\geq1).
\end{equation}

\section{Recurrent ideals and finite normal forms}\label{sec:normal}

Throughout this section, $X$ is a nonempty finite poset. By an \emph{ideal} we mean a downward-closed subset, without any directedness assumption. Thus
\[
 \J(X)=\{D\subseteq X:(b\in D\ \&\ a\leq_X b)\Longrightarrow a\in D\}.
\]

For a length-$\omega$ word $u$, let
\[
 \rec(u)=\{a\in X:a\text{ occurs infinitely often in }u\},
 \qquad
 I(u)=\down\rec(u).
\]
The recurrent set is nonempty because the alphabet is finite. Its downward closure, rather than the literal recurrent set, is the correct invariant when labels may increase under embeddings.

For each nonempty ideal $D$, fix an ordering $d_0,\ldots,d_{t-1}$ of its maximal elements, and define
\begin{equation}\label{eq:cD}
 c_D=(d_0d_1\cdots d_{t-1})^\omega.
\end{equation}
Different choices of ordering give equivalent words. Every nonempty tail of $c_D$ is equivalent to $c_D$.

\begin{lemma}[Periodic support]\label{lem:periodic}
For nonempty ideals $D,E$, one has
\begin{equation}\label{eq:periodic}
 c_D\sub c_E\quad\Longleftrightarrow\quad D\subseteq E.
\end{equation}
A letter $a$ embeds into every tail of $c_D$ if and only if $a\in D$. Consequently, if every letter of a finite word $p$ lies in $D$, then $p c_D\equ c_D$.
\end{lemma}
\begin{proof}
If $D\subseteq E$, each maximal element of $D$ is below some maximal element of $E$. The latter element appears infinitely often in $c_E$. Match source letters successively, each time choosing a suitable target occurrence later than all previously chosen occurrences. This constructs the embedding, and the same argument works after discarding any finite target prefix.

Conversely, if $c_D\sub c_E$, every maximal element of $D$ is below some target label in $\Max E$, hence belongs to $E$. Downward closure gives $D\subseteq E$. The one-letter assertion follows by the same reasoning. To absorb $p$, place its finitely many letters in an initial portion of $c_D$ and embed the remaining copy of $c_D$ into a tail; the reverse embedding simply omits $p$.
\end{proof}

\begin{lemma}[One-block normal form]\label{lem:oneblock}
Every length-$\omega$ word $u$ is equivalent to $p c_D$, where $D=I(u)$ and $p$ is a finite word. We may require that $p$ be empty or end in a letter outside $D$.
\end{lemma}
\begin{proof}
Every letter outside $D$ occurs only finitely often. Since there are finitely many letters, there is a last occurrence outside $D$, if there is any such occurrence. Let $p$ be the prefix ending at that last occurrence, or let $p=\emp$ if there is none. Write the remaining tail as $t$.

Every label of $t$ lies in $D$. Every maximal element $d$ of $D$ occurs infinitely often in $t$: by $D=\down\rec(u)$, there is a recurrent $b\geq_X d$, and maximality forces $b=d$. Thus $t\sub c_D$ by greedy matching to dominating maximal labels, and $c_D\sub t$ by matching each of its labels to a later occurrence of the same recurrent maximal label. Concatenating with $p$ gives the assertion.
\end{proof}

Every ordinal below $\omega^2$ has a unique form $\omega r+k$ with $r,k\in\N$. Splitting a word at the canonical block boundaries and applying \cref{lem:oneblock} gives an expression
\begin{equation}\label{eq:normal}
 p_0c_{D_0}p_1c_{D_1}\cdots p_{r-1}c_{D_{r-1}}p_r,
\end{equation}
where every $p_i$ is finite and each $p_i$ with $i<r$ is empty or ends outside $D_i$. The final finite word $p_r$ has length $k$.

The construction is existential for an arbitrary infinite input word: it uses knowledge of which letters recur infinitely often. The algorithm later in the paper takes finite expressions of the form \eqref{eq:normal} or unnormalized token expressions as input; it is not an algorithm for discovering recurrence from an arbitrary stream.

\section{Allowed supports and the generator upper bound}\label{sec:upper}

\begin{definition}[A support family]
Let $\F\subseteq\J(X)\setminus\{\varnothing\}$. Define $\W_{\F}(X)$ to be the quasiorder of finite concatenations of individual letters of $X$ and the words $c_D$ with $D\in\F$, with the inherited transfinite subsequence relation. Equivalent expressions represent the same quotient element. We suppress $X$ when it is fixed.
\end{definition}

In intrinsic terms, these are the words whose canonical $\omega$-blocks have recurrent ideals in $\F$, up to equivalence. Finite prefixes cannot change the recurrent ideal of a block, and \cref{lem:oneblock} proves the converse. In particular,
\begin{equation}\label{eq:allfamily}
 \W_{\varnothing}=X^*,
 \qquad
 \W_{\J(X)\setminus\{\varnothing\}}\equiv s_{\omega^2}(X)
\end{equation}
as quotient orders.

\begin{proposition}[Generator bound]\label{prop:upper}
If $n=|X|\geq1$ and $t=|\F|$, then $\W_{\F}$ is a wqo and
\begin{equation}\label{eq:upper}
 \ot(\W_{\F})\leq\HH_{n+t}.
\end{equation}
\end{proposition}
\begin{proof}
Take a discrete alphabet of $n+t$ formal tokens: one token for each letter of $X$ and one for each $D\in\F$. Interpret a letter token as its one-letter word and an ideal token as $c_D$, then concatenate. Deleting tokens from a source token word gives a subsequence embedding of its interpretation into the target interpretation. Hence interpretation is monotone from the finite Higman order on the discrete token alphabet.

It is surjective onto $\W_{\F}$ up to equivalence, by definition. A monotone image of a wqo is a wqo. Monotone-surjection monotonicity and \eqref{eq:higman} give \eqref{eq:upper}.
\end{proof}

The map need not be order-reflecting. For example, a finite word with letters in $D$ can be absorbed by a subsequent $c_D$. Counting distinct generators therefore establishes only an upper bound. The next sections construct the missing lower bound while explicitly preventing such absorption.

\section{Cofinality forces the separators to align}\label{sec:cofinal}

\begin{lemma}[Cofinal target block]\label{lem:cofinal}
If $c_D$ embeds into a word $v$ of length below $\omega^2$, a tail of its image is cofinal in one canonical target $\omega$-block. If that block has recurrent ideal $E$, then $D\subseteq E$.
\end{lemma}
\begin{proof}
Write $\len(v)=\omega r+k$. Under a strictly increasing map from $\omega$, the index of the target canonical block is a nondecreasing sequence in a finite set. It is eventually constant. The eventual block cannot be the final finite tail, because the image is infinite. Within an $\omega$-block, an infinite increasing set of positions is cofinal.

Fix $d\in\Max D$. Infinitely many occurrences of $d$ in the source are mapped into that eventual block. The alphabet is finite, so some target label $b\geq_X d$ occurs infinitely often among these images. It is recurrent in the block, and hence $d\in E$. This holds for every maximal element of $D$, proving $D\subseteq E$.
\end{proof}

If one source periodic block is followed by another, their images cannot both be cofinal in the same target $\omega$-block. All points of the second source block must follow all points of the first, whereas a cofinal subset leaves no later point inside that target block.

\begin{lemma}[Synchronization]\label{lem:sync}
Suppose $D\notin\F$ is a nonempty ideal satisfying
\begin{equation}\label{eq:separation}
 D\nsubseteq E\qquad\text{for every }E\in\F.
\end{equation}
Let $b_0,\ldots,b_m$ and $d_0,\ldots,d_m$ belong to $\W_{\F}$. In any embedding
\begin{equation}\label{eq:syncexpr}
 b_0c_D b_1c_D\cdots c_D b_m
 \sub
 d_0c_D d_1c_D\cdots c_D d_m,
\end{equation}
each of the $m$ displayed source separators is cofinally mapped into the corresponding displayed target separator block, in order.
\end{lemma}
\begin{proof}
By \cref{lem:cofinal}, every source separator must eventually use a target $\omega$-block with recurrent ideal containing $D$. Condition \eqref{eq:separation} excludes all blocks supplied by the old words $d_i$. Thus it must use one of the $m$ displayed target copies of $c_D$.

Successive source separators must use different target blocks, and their assignment is strictly increasing. An increasing injection from an $m$-element chain into itself is the identity. This gives the required alignment.
\end{proof}

\paragraph{The role of explicit segments.}
A displayed $c_D$ can follow a finite suffix of $d_i$. Its pure periodic segment then starts at a finite offset inside the corresponding canonical $\omega$-block. The segment still has length $\omega$ and is cofinal in that block. Its end is the next canonical boundary. In the arguments below, ``inside the separator'' means inside this \emph{explicit pure periodic segment}, not inside the finite prefix that precedes it. This distinction accounts for the ordinal identity $k+\omega=\omega$.

\begin{corollary}[Location between synchronized boundaries]\label{cor:between}
Under \eqref{eq:syncexpr}, all images of $b_i$ occur after the end of the previous target separator, when $i>0$. If $i<m$, they occur before the image of the first point of the next source separator and hence before the end of the next target separator. The final $b_m$ embeds into the final target word $d_m$.
\end{corollary}
\begin{proof}
Every point of $b_i$ follows every point of the preceding source separator. The preceding image is cofinal in the corresponding target block, so such a point must occur at or after its end. Every point of $b_i$ also precedes the next source separator, if present. Since that separator is cofinal in the corresponding target block, its first image precedes the end of that block. After the last target separator, the only remaining target segment is $d_m$.
\end{proof}

This still permits a finite tail of $b_i$ to enter the next target separator. Guards eliminate precisely that possibility.

\section{Two guards against absorption}\label{sec:guards}

\subsection{A proper ideal: use a letter outside it}

\begin{lemma}[Terminal-letter guard]\label{lem:letterguard}
For any $a\in X$, the map $u\mapsto ua$ is order-preserving and order-reflecting on words below $\omega^2$. If $D$ is proper and $a\notin D$, the last letter of $ua$ cannot embed into $c_D$.
\end{lemma}
\begin{proof}
An embedding $u\sub v$ extends to $ua\sub va$ by matching the two new terminal letters. Conversely, suppose $ua\sub va$. The distinguished last source letter has some image no later than the last target position. All points of $u$ precede this image, so they lie in the target prefix $v$. Restricting the embedding gives $u\sub v$.

If $a\notin D$, no label of $c_D$ can dominate $a$, since all labels of $c_D$ lie in the downward-closed set $D$.
\end{proof}

\subsection{The full ideal: use a limit-length guard}

A full-support block $c_X$ can absorb every finite word, so no terminal letter is outside its support. It is here that a second construction is needed.

\begin{lemma}[Limit guard]\label{lem:limitguard}
Suppose $\F$ contains a proper nonempty ideal $R$. Choose $a\in X\setminus R$ and put
\begin{equation}\label{eq:limitguard}
 g(u)=u a c_R\qquad(u\in\W_{\F}).
\end{equation}
Then $g:\W_{\F}\to\W_{\F}$ is order-preserving and order-reflecting. Every $g(u)$ has nonzero limit length, and every tail starting at a position of $g(u)$ is infinite.
\end{lemma}
\begin{proof}
Order preservation follows by concatenation. Suppose $uac_R\sub vac_R$. The distinguished source $a$ cannot map into the final target $c_R$, because $a\notin R$. It therefore maps within $va$, and all source positions belonging to $u$ map into $v$. Thus $u\sub v$.

If $\len(u)=\omega r+k$, then
\[
 \len(g(u))=(\omega r+k)+1+\omega=\omega(r+1).
\]
This is a positive limit ordinal. More concretely, its terminal segment is a full copy of $\omega$, so infinitely many positions remain after any selected position.
\end{proof}

\begin{lemma}[A limit zone cannot leak into a separator]\label{lem:noleak}
In the situation of \cref{lem:sync}, suppose $b_i$ has nonzero limit length and $i<m$. Then the image of $b_i$ cannot meet the explicit pure target separator $c_D$ immediately following $d_i$.
\end{lemma}
\begin{proof}
Suppose a source position $\xi$ in $b_i$ maps to a position $\eta$ inside that pure target separator. Let $\zeta$ be the image of the first position of the next source separator. By strict monotonicity, $\eta<\zeta$; by synchronization, $\zeta$ occurs before the end of this same target separator. Thus both positions belong to a segment of order type $\omega$, and there are only finitely many target positions strictly between them.

Since $b_i$ has nonzero limit length, infinitely many source positions lie after $\xi$ but before the first point of the next source separator. Their images would have to lie strictly between $\eta$ and $\zeta$, a contradiction.
\end{proof}

The lemma is not an assertion that a limit-length word never embeds into a universal periodic word. It certainly may. The obstruction is the requirement that another synchronized source separator must begin at an actual target position before the end of that same target block.

\section{The guarded product embedding}\label{sec:product}

\begin{theorem}[Product construction]\label{thm:product}
Let $D\notin\F$ satisfy \eqref{eq:separation}. Assume either
\begin{enumerate}[label=(\roman*)]
 \item $D\neq X$, or
 \item $D=X$ and $\F$ contains a proper nonempty ideal.
\end{enumerate}
Then, for every $r\geq1$, there is an order embedding
\begin{equation}\label{eq:productembedding}
 \W_{\F}^{\,r}\lhook\joinrel\longrightarrow\W_{\F\cup\{D\}},
\end{equation}
where the source has componentwise order.
\end{theorem}
\begin{proof}
If $D$ is proper, choose $a\notin D$ and set $g(u)=ua$. In case (ii), choose $R\in\F$ proper and $a\notin R$, and set $g(u)=uac_R$. Both maps take old-family words to old-family words and preserve and reflect the order, by \cref{lem:letterguard,lem:limitguard}.

For $r=m+1$, define
\begin{equation}\label{eq:Phi}
 \Phi_m(u_0,\ldots,u_m)
  =g(u_0)c_D g(u_1)c_D\cdots c_D g(u_m).
\end{equation}
If $u_i\sub v_i$ for every $i$, concatenate the guarded component embeddings and identity embeddings on the separator blocks. This proves that $\Phi_m$ is order-preserving.

For reflection, suppose
\[
 \Phi_m(u_0,\ldots,u_m)\sub\Phi_m(v_0,\ldots,v_m).
\]
The case $m=0$ is the order reflection of $g$. For $m>0$, \cref{lem:sync} aligns all separators, and \cref{cor:between} places the image of each guarded component after the preceding target separator.

Consider a nonfinal component $g(u_i)$. In case (i), its final letter $a$ cannot lie inside the next pure target separator, because $a\notin D$. It also cannot lie after that separator, because the next source separator must be cofinally mapped into it. Therefore this final letter, and all preceding positions of the guarded component, lie in the corresponding target component $g(v_i)$. Hence $g(u_i)\sub g(v_i)$.

In case (ii), $g(u_i)$ has nonzero limit length. By \cref{lem:noleak}, it cannot meet the next pure target separator at all. Together with the two boundary restrictions, this again gives $g(u_i)\sub g(v_i)$. The last guarded component embeds into the last target component directly by \cref{cor:between}.

Finally, order reflection of $g$ gives $u_i\sub v_i$ for every $i$. Thus $\Phi_m$ is order-reflecting, and descends to an injective order embedding on equivalence classes.
\end{proof}

\begin{lemma}[Finite powers amplify a maximal order type]\label{lem:power}
For any wqo $P$ and any finite $r\geq1$,
\begin{equation}\label{eq:lexlower}
 \ot(P^r)\geq\ot(P)^r,
\end{equation}
with ordinary ordinal exponentiation on the right.
\end{lemma}
\begin{proof}
Let $L$ be a maximal linear extension of $P/{\equ}$, of type $\alpha=\ot(P)$. Order $r$-tuples by comparing their last unequal coordinate, using $L$ in that coordinate. This is a well-order of type $\alpha^r$: the last coordinate indexes $\alpha$ consecutive copies of the corresponding $(r-1)$-coordinate order.

It extends the componentwise order. Indeed, if one tuple is coordinatewise below another, all its coordinates are below or equal in $L$; at the last coordinate where they differ, the first tuple is strictly earlier. This supplies the claimed linear extension of $P^r$.
\end{proof}

\begin{corollary}[One support raises one Higman level]\label{cor:activation}
Under the hypotheses of \cref{thm:product}, if $\ot(\W_{\F})=\HH_q$ and $q=n+|\F|$, then
\[
 \ot(\W_{\F\cup\{D\}})=\HH_{q+1}.
\]
\end{corollary}
\begin{proof}
The product embedding and \cref{lem:power} give the lower bound $\HH_q^r$ for every $r\geq1$. Taking the supremum and applying \eqref{eq:amplification} gives $\HH_{q+1}$. The upper bound is \cref{prop:upper} with one more token.
\end{proof}

Only an explicit linear extension of a Cartesian product was used here. In particular, the proof does not require any general formula for maximal order types of lexicographic products of arbitrary partial orders.

\section{The full support-family classification}\label{sec:classification}

\begin{theorem}[Classification by allowed recurrent ideals]\label{thm:families}
Let $X$ be a nonempty finite poset of cardinality $n$, and let $\F\subseteq\J(X)\setminus\{\varnothing\}$. Then
\begin{equation}\label{eq:families}
 \ot(\W_{\F})=
 \begin{cases}
   \HH_n\cdot\omega,&\F=\{X\},\\
   \HH_{n+|\F|},&\F\neq\{X\}.
 \end{cases}
\end{equation}
\end{theorem}
\begin{proof}[Proof outside the exceptional family]
For $\F=\varnothing$, the result is \eqref{eq:higman}. Otherwise, enumerate the ideals in $\F$ in nondecreasing cardinality, placing $X$ last when it occurs. At a stage inserting $D$, every old ideal $E$ has $|E|\leq|D|$ and is distinct from $D$. Therefore $D\nsubseteq E$.

For every proper $D$, \cref{cor:activation} applies with the terminal-letter guard. If $X$ occurs, the assumption $\F\neq\{X\}$ ensures that some proper nonempty ideal was inserted earlier, so the limit guard applies at the final stage. Beginning at $\HH_n$, each insertion raises the subscript by one. This gives $\HH_{n+|\F|}$.
\end{proof}

\begin{proof}[Proof for the exceptional family]
Assume $\F=\{X\}$. Every finite prefix before a copy of $c_X$ is absorbed by that copy. Thus each word is equivalent to
\begin{equation}\label{eq:universalnormal}
 c_X^r u,\qquad r\in\N,\quad u\in X^*.
\end{equation}
Here $c_X^r$ means a concatenation of $r$ copies, with $c_X^0=\emp$.

If $r<s$, then $c_X^r u\sub c_X^s v$ for every finite $u,v$: match the first $r$ copies cofinally, then place $u$ in a further universal block. If $r>s$, an embedding is impossible because the source has more $\omega$-blocks. For $r=s$, each source block must use the corresponding target block cofinally, and the finite source tail must enter the finite target tail. Hence
\begin{equation}\label{eq:universalorder}
 c_X^r u\sub c_X^s v
 \quad\Longleftrightarrow\quad
 r<s\ \text{or}\ (r=s\ \text{and}\ u\sub v).
\end{equation}

The quotient is therefore an ordinal sum of $\omega$ copies of $X^*$. Any linear extension respects all layer boundaries, so its type is at most $\HH_n\cdot\omega$. Choosing a maximal extension in every layer attains that type. This proves the exceptional line of \eqref{eq:families}.
\end{proof}

The exceptional value is genuinely smaller than the naive generator bound:
\[
 \HH_n\cdot\omega
 =\omega^{\omega^{n-1}+1}
 <\omega^{\omega^n}=\HH_{n+1}.
\]
The missing proper recurrent support is exactly what prevents construction of the limit guard.

\subsection{Deduction of the finite-alphabet formula}

\begin{proof}[Proof of \cref{thm:main}]
If $X=\varnothing$, only the empty word exists, giving type $1$.

If $|X|=1$, there is exactly one word of each length $\alpha<\omega^2$. Such a word embeds into another exactly when its length is at most the other's length. The quotient is the chain $\omega^2$.

Now suppose $n\geq2$. A minimal element $a\in X$ exists, and $\{a\}$ is a proper nonempty ideal. Hence the full family $\F=\J(X)\setminus\{\varnothing\}$ is not the exceptional singleton family. Since $|\F|=j-1$, \cref{thm:families} and \eqref{eq:allfamily} give
\[
 \ot\bigl(s_{\omega^2}(X)\bigr)
  =\HH_{n+j-1}
  =\omega^{\omega^{n+j-2}}.
\]
\end{proof}

\begin{corollary}[Finite quasiorders]
For a finite quasiorder, first quotient the alphabet by mutual comparability. The formula of \cref{thm:main} then applies to that finite poset. In particular, $n$ must count alphabet equivalence classes, not necessarily the original number of labels.
\end{corollary}
\begin{proof}
Passing to label classes preserves and reflects the existence of a strictly increasing label-monotone embedding. The induced word quotients are therefore isomorphic.
\end{proof}

\section{Worked examples and why the safeguards matter}\label{sec:examples}

\subsection{The two-letter antichain}

Take $X=\{0,1\}$ with incomparable elements. Its three nonempty ideals yield
\[
 c_{\{0\}}=0^\omega,
 \qquad c_{\{1\}}=1^\omega,
 \qquad c_X=(01)^\omega.
\]
The lower-bound induction can be read directly from the following table.

\begin{center}
\begin{tabular}{@{}lll@{}}
\toprule
Allowed recurrent ideals & Guard for the new ideal & Maximal type\\
\midrule
$\varnothing$ & None; finite words & $\HH_2$\\
$\{\{0\}\}$ & $u\mapsto u1$ & $\HH_3$\\
$\{\{0\},\{1\}\}$ & $u\mapsto u0$ & $\HH_4$\\
$\{\{0\},\{1\},X\}$ & $u\mapsto u1\,0^\omega$ & $\HH_5$\\
\bottomrule
\end{tabular}
\end{center}

At each row, arbitrary finite tuples from the preceding row are separated by the newly admitted periodic block. The guard makes the tuple map order-reflecting. The last row is all words below $\omega^2$, and gives $\HH_5=\omega^{\omega^4}$.

The proof is not a claim that the naive interpretation from the five-token Higman order is injective. That interpretation is not injective even on equivalence classes. Instead, it establishes the same maximal order type through an iterated family of finite-product lower bounds.

\subsection{The two-letter chain}

Now take $0<1$. There are only two nonempty ideals, $\{0\}$ and $X$. A full-support periodic representative is $c_X=1^\omega$: the lower label $0$ can always be matched to $1$. Starting from $\HH_2$, insert $\{0\}$ using $u1$, then insert $X$ using $u1\,0^\omega$. The result is $\HH_4=\omega^{\omega^3}$.

Thus two finite alphabets can have the same cardinality and the same finite-alphabet maximal type, yet their transfinite-word orders have different maximal types. The number of ideals measures the difference at this cutoff.

\subsection{A small table of the general formulas}

\begin{center}
\begin{tabular}{@{}rrrrr@{}}
\toprule
$n$ & $|\J(C_n)|$ & Chain exponent $2n-1$ & $|\J(A_n)|$ & Antichain exponent $n+2^n-2$\\
\midrule
2 & 3 & 3 & 4 & 4\\
3 & 4 & 5 & 8 & 9\\
4 & 5 & 7 & 16 & 18\\
5 & 6 & 9 & 32 & 35\\
\bottomrule
\end{tabular}
\end{center}

An ``exponent'' $e$ in this table means that the answer is $\omega^{\omega^e}$, not $\omega^e$. The entries follow arithmetically from \eqref{eq:anti} and \eqref{eq:chain}.

\subsection{Three invalid shortcuts}

\paragraph{Omitting guards.}
For any letter $a$, one has
\[
 a c_X\equ c_X,
 \qquad a\nsub\emp.
\]
Thus a component placed directly before a separator can disappear under embeddability, so a tuple map without guards need not reflect the componentwise order.

\paragraph{Using a finite guard for full support.}
The same problem persists after appending one letter:
\[
 aa c_X\sub a c_X,
 \qquad a\nsub\emp.
\]
More generally, any fixed finite padding can be absorbed into a universal periodic tail. The guard in \eqref{eq:limitguard} is infinite for a mathematical reason.

\paragraph{Inserting a support after one that contains it.}
Condition \eqref{eq:separation} cannot simply be discarded. On the two-letter antichain, let the old family contain $E=X$, and let the new support be $D=\{0\}$. With the naive terminal guard $1$, consider the two pairs
\[
 (\emp,0),\qquad(c_E,\emp).
\]
The first pair is not componentwise below the second. Nevertheless, their guarded encodings satisfy
\begin{equation}\label{eq:negative}
 1\,0^\omega\,0\,1
 \sub
 c_E\,1\,0^\omega\,1.
\end{equation}
Match the initial $1$ and the source $0^\omega$ inside $c_E$, using the latter cofinally. Then match the remaining $0$ into the displayed target $0^\omega$ and the last $1$ to the terminal $1$. The source separator has used the old containing support, rather than its designated target separator. Ordering insertions by cardinality prevents this failure.

\section{Exact-length strata}\label{sec:exact}

This section gives supplementary formulas and explains why the maximal type of all words cannot be recovered by taking a supremum of the types at individual lengths. Let
\[
 \E_\alpha(X)=\{u:\len(u)=\alpha\},
 \qquad n=|X|\geq1,
 \qquad d=|\J(X)|-2.
\]
Thus $d$ is the number of proper nonempty recurrent ideals.

\subsection{A single infinite block}

\begin{proposition}\label{prop:exactomega}
For every nonempty finite poset,
\begin{align}
 \ot(\E_\omega(X))&=\HH_n\cdot d+1,\label{eq:exactomega}\\
 \ot(s_{\omega+1}(X))&=\HH_n\cdot(d+1)+1.\label{eq:atmostomega}
\end{align}
\end{proposition}
\begin{proof}
Fix a proper nonempty ideal $D$. In its fibre, use the forms $p c_D$ with $p$ empty or ending outside $D$. For two such forms,
\begin{equation}\label{eq:fibre}
 p c_D\sub q c_D\quad\Longleftrightarrow\quad p\sub q.
\end{equation}
The forward direction is immediate when $p=\emp$. Otherwise its final letter lies outside $D$, so it cannot enter the target periodic tail; therefore all of $p$ embeds into $q$. The reverse direction follows by concatenation.

The fibre consequently has type at most $\HH_n$. Choose $a\notin D$ and map $u\in X^*$ to $ua c_D$. By \eqref{eq:fibre} and terminal-letter cancellation, this is order-preserving and order-reflecting. Hence the fibre has type exactly $\HH_n$.

The full-support fibre is a single class $c_X$, which is above every length-$\omega$ word. Moreover, an embedding between two fibres forces inclusion of their recurrent ideals by \cref{lem:cofinal}. Order the $d$ proper ideals by nondecreasing cardinality and concatenate maximal extensions of their fibres, then place $c_X$ last. No backwards comparison is violated. This constructs a linear extension of type $\HH_n d+1$.

For the upper bound, erase all comparisons between fibres. This gives a disjoint union of $d$ orders of type $\HH_n$ and one singleton, whose type is $\HH_n d+1$ by natural addition. Adding comparisons cannot increase maximal order type. This proves \eqref{eq:exactomega}.

For \eqref{eq:atmostomega}, add the finite-word part first. No infinite word embeds into a finite word, so this placement respects the order and adds a layer of type $\HH_n$. The matching upper bound follows from the disjoint union of this additional part with the previous fibres.
\end{proof}

\subsection{Several blocks and a fixed finite tail}

\begin{lemma}[Rigidity at equal block count]\label{lem:equalblocks}
If two words each have $r$ canonical $\omega$-blocks, an embedding between them sends each source block cofinally into the target block with the same index. The entire source block is mapped within that target block. The source finite tail maps within the target finite tail.
\end{lemma}
\begin{proof}
Apply the cofinal-block argument to each whole source block, or to its periodic normal-form tail. Distinct source blocks require distinct target blocks, in increasing order; there are $r$ of each, so the assignment is the identity.

A source block cannot start before the end of the previous target block, because the preceding source block has already used that block cofinally. It cannot have an image point after its own assigned target block, because then its remaining infinite tail could not be cofinal in the earlier assigned block. Hence the full source block stays inside its assigned target block. After the last block, only the target finite tail remains.
\end{proof}

\begin{proposition}[Fixed-length product]\label{prop:exactproduct}
For $r,k\in\N$,
\begin{equation}\label{eq:exactisom}
 \E_{\omega r+k}(X)\equiv\E_\omega(X)^r\times X^k
\end{equation}
as quotient orders. Consequently, with $A=\HH_n d+1$,
\begin{equation}\label{eq:exacttype}
 \ot(\E_{\omega r+k}(X))=A^{\otimes r}\otimes n^k.
\end{equation}
\end{proposition}
\begin{proof}
Concatenation gives a bijection between actual words of the stated lengths and their canonical blocks and tail. By \cref{lem:equalblocks}, it reflects and preserves the blockwise embedding relation. An increasing injection between finite chains of the same length $k$ is the identity, so the relation on their label tuples is exactly the componentwise relation on $X^k$.

This proves \eqref{eq:exactisom}. Its finite tail has $n^k$ elements. Apply the natural-product identity \eqref{eq:naturalops} and \cref{prop:exactomega} to obtain \eqref{eq:exacttype}.
\end{proof}

Because $\HH_n=\omega^{\omega^{n-1}}$, natural multiplication expands \eqref{eq:exacttype} to the Cantor normal form
\begin{equation}\label{eq:cnfexact}
 \ot(\E_{\omega r+k}(X))
  =\sum_{i=r,r-1,\ldots,0}
    \omega^{\omega^{n-1}\cdot i}
       \left(n^k\binom ri d^i\right).
\end{equation}
Terms with zero coefficient are omitted, and $d^0=1$, including when $d=0$. The sum is written in decreasing $i$, so it is an ordinary Cantor-normal-form sum. Indeed, in the expansion of $(\HH_n d+1)^{\otimes r}$, selecting the leading monomial in $i$ factors contributes coefficient $\binom ri d^i$ and exponent equal to the natural sum of $i$ copies of $\omega^{n-1}$, namely $\omega^{n-1}i$.

\paragraph{A warning about taking suprema over lengths.}
For $n\geq2$, one has $d\geq1$, and \eqref{eq:cnfexact} implies
\begin{equation}\label{eq:stratumsup}
 \sup_{r,k<\omega}\ot(\E_{\omega r+k}(X))
   =\omega^{\omega^n}.
\end{equation}
The equality follows by taking the supremum of the leading exponents $\omega^{n-1}r$; finite coefficients and lower terms do not change that supremum. In contrast, the type of \emph{all} words is $\omega^{\omega^{n+j-2}}$, which is strictly larger. To see strictness, take a linear extension of $X$: its $n+1$ prefixes are distinct ideals, so $j\geq n+1$ and $n+j-2>n$ when $n\geq2$.

There is no contradiction. The length strata are not ordinal-sum layers: a shorter word need not embed into every longer word. Maximal linear extensions can exploit incomparabilities between different strata. The support induction detects this additional complexity, whereas the separate exact-length calculations do not.

\section{An exact symbolic embedding algorithm}\label{sec:algorithm}

\subsection{Finite representations}

A token expression is a finite list of two types of tokens: a letter $a\in X$, or an ideal token $[D]$ denoting $c_D$. This finite syntax describes actual transfinite words. The algorithm does not replace $\omega$ by a large integer.

Maintain a cursor in the target token list and process source tokens in order. A successful finite match inside a target periodic token leaves the cursor at that token, since only a finite initial portion has been used. A successful source periodic token consumes the whole matched target periodic token cofinally.

\begin{center}
\begin{tabular}{@{}llll@{}}
\toprule
Source & Target & Matching condition & Cursor after a match\\
\midrule
$a$ & $b$ & $a\leq_X b$ & Advance past $b$\\
$a$ & $[E]$ & $a\in E$ & Remain at $[E]$\\
$[D]$ & $[E]$ & $D\subseteq E$ & Advance past $[E]$\\
$[D]$ & $b$ & Never & No match\\
\bottomrule
\end{tabular}
\end{center}

When the current target token cannot match, skip it. Failure occurs if the target is exhausted while a source token remains.

\begin{theorem}[Correctness of the greedy token algorithm]\label{thm:algorithm}
The algorithm returns true exactly when the interpreted source word embeds into the interpreted target word. For a fixed alphabet and unit-cost support tests, its time complexity is $O(N+M)$ for source and target token counts $N,M$, and its auxiliary space is $O(1)$.
\end{theorem}
\begin{proof}
First consider one source letter $a$. In a finite target token, the condition is exactly $a\leq_X b$. In a periodic target token, $a$ can be matched if and only if $a\in E$, by \cref{lem:periodic}. Any tail remaining after a finite number of such matches is equivalent to $c_E$. Matching at the first suitable target token never removes an embedding opportunity that matching later would preserve: an earlier position leaves at least the same later target segments available.

Now consider a source token $[D]$. By \cref{lem:cofinal}, any embedding of it into the remaining target must eventually be cofinal in a target periodic token with support containing $D$. Conversely, $c_D$ embeds into every tail of the first such token. Using that earliest possible token cofinally leaves all subsequent tokens available, a suffix no shorter than one obtained by using a later possible target token. Target tokens preceding it can therefore be skipped without loss.

These two observations give an induction on the finite number of source tokens. After a letter match, the remaining target is exactly the next token suffix, possibly with a finite prefix removed from its leading periodic token; that removal does not change its equivalence class. After a periodic match, the remaining target begins after the cofinally consumed token. Thus each greedy step preserves existence of an embedding for the remaining source, proving necessity. Concatenating the indicated finite and cofinal matches proves sufficiency.

Every successful match advances the source cursor once. Every unsuccessful examination advances the target cursor, which never moves backwards. These two counts bound the total number of operations by a constant times $N+M$. Apart from input lists, only cursors and token values are stored.
\end{proof}

When the alphabet is part of the input, set-membership and inclusion costs must also be counted. The implementation uses bit masks; an inclusion test takes time proportional to the machine-word length of the masks. The linear token bound does not claim constant-cost manipulation of arbitrarily large integers. Nor does it include enumerating all ideals of a large poset.

\subsection{Canonical quotient representatives}

In each finite prefix preceding $[D]$, remove the longest suffix all of whose letters lie in $D$. This is a stack-based normalization: when an ideal token arrives, pop the immediately preceding letters while they belong to its ideal, then append the ideal token. Do not alter the final finite tail.

\begin{proposition}[Uniqueness of normalized expressions]\label{prop:canonical}
The resulting expression is equivalent to the input. Two normalized token expressions over a finite poset are mutually embeddable if and only if they are identical.
\end{proposition}
\begin{proof}
Every removed suffix is absorbed by the following periodic token by \cref{lem:periodic}, so normalization preserves equivalence.

Mutual strict embeddings imply equality of ordinal lengths. Therefore the two expressions have the same number of $\omega$-blocks and the same finite tail length. By \cref{lem:equalblocks}, corresponding whole blocks embed mutually. Their recurrent ideals must contain each other, so the ideal tokens agree. In the two finite prefixes before such a shared ideal, the last letter, if present, lies outside that ideal. Exactly as in \eqref{eq:fibre}, mutual block embeddings force mutual embeddings of the finite prefixes.

Two finite words over a poset that embed mutually have the same finite length. Both index maps are therefore the identity, and antisymmetry forces equality of corresponding labels. This proves equality of each finite prefix and of the final tail. The converse is immediate.
\end{proof}

The quotient below $\omega^2$ consequently has a countable and effectively comparable system of finite representatives, even though there are uncountably many literal infinite words on an alphabet with at least two elements.

\subsection{Implementation structure}

The module \texttt{ordinal\_words.py} contains a validated finite-poset representation, ideal enumeration, an iterative greedy comparator, a separate memoized branching comparator, normalization, the two guards, and the product construction. The branching comparator explores both matching and skipping transitions rather than committing to the earliest match. It has $O(NM)$ states and is used as a testing oracle on finite token lists.

Both comparators rely on the proved periodic-token abstraction. Their agreement checks the implementation, not the correctness of that abstraction in isolation. \cref{lem:cofinal,thm:algorithm} supply the mathematical justification.

\section{Reproducible computational audit}\label{sec:verification}

The accompanying program \texttt{verify.py} uses exact integer and finite-set operations. All tests reported here were run; the machine-readable output is \texttt{verification\_results.json}. The pseudorandom seed is $20260919$.

\begin{center}
\begin{tabular}{@{}p{0.63\linewidth}r@{}}
\toprule
Test family & Ordered pairs tested\\
\midrule
Two-letter antichain: all token expressions of length at most $4$ & 609,961\\
Two-letter chain: all token expressions of length at most $4$ & 116,281\\
Exhaustive small tuples across binary support-family insertions & 44,667\\
Seeded guarded-product tests across small posets & 88,750\\
Additional small-poset embedding and normalization tests & 12,500\\
Universal-only support-family order formula & 7,500\\
Fixed-length product decomposition & 12,500\\
\midrule
Total pair test cases & 892,159\\
\bottomrule
\end{tabular}
\end{center}

All tests passed. Three separate negative controls detected the failures described in \cref{sec:examples}.

\paragraph{The exhaustive binary domain.}
The antichain has five tokens, giving $1+5+5^2+5^3+5^4=781$ expressions and $781^2=609{,}961$ ordered pairs. The chain has four tokens, giving $341$ expressions and $116{,}281$ pairs. For every pair, the greedy and branching comparators agree. Normalization preserves the result, and mutually embeddable expressions have identical normal forms.

\paragraph{The product tests.}
Across the two binary alphabets, all allowed support subfamilies are considered, with the universal-only exceptional activation treated separately. There are $14$ tested activation occurrences. At each, component words have at most one old token, and all ordered pairs of tuples of arity $1$, $2$, or $3$ are checked. Repeated support prefixes in different subfamilies are intentionally included.

A second suite enumerates all distinct naturally labelled posets on $n=1,2,3,4$ elements. Here ``naturally labelled'' means $a<_X b$ implies $a<b$ numerically. The computed counts are $1,2,7,40$, respectively. Every isomorphism type of a poset with at most four elements has such a labelling, although isomorphic labelled copies are not removed. The resulting $50$ posets and their ideal counts are recorded in \texttt{finite\_posets.csv}.

Their all-support insertion chains have $355$ admissible activation stages. Each is tested on $250$ seeded tuple pairs of arity at most $4$. Source components have at most $10$ tokens; some guaranteed-positive cases append up to $4$ additional tokens to obtain target components. The comparator result is checked against the conjunction of all component comparisons, which is precisely the assertion of \cref{thm:product} for that tuple pair.

\paragraph{What these checks do not establish.}
No finite test counts maximal order types or verifies a transfinite supremum. No test replaces a proof that all separators synchronize. The tests check finite representations of the constructions, including cases of absorption and component leakage that are easy to mishandle. The unrestricted conclusions come from the lemmas and ordinal calculations in the preceding sections.

\paragraph{Reproduction.}
The code requires Python 3.10 or later and no third-party packages. From the archive directory, run
\begin{lstlisting}
python verify.py
\end{lstlisting}
The default run recreates both the JSON result file and the CSV table. To change the seeded sampling while keeping the exhaustive binary domain, use, for example,
\begin{lstlisting}
python verify.py --seed 17 --samples 1000
\end{lstlisting}
To compile this article, run \texttt{pdflatex article.tex} three times, or use the included \texttt{build.sh}. The supplied PDF is already compiled.

\section{Proof audit and limits of the conclusion}\label{sec:audit}

\subsection{The logical dependency chain}

The main theorem depends on three classical inputs: Higman's wqo theorem, the finite-word maximal-type formula, and maximal-linear-extension existence with its monotonicity consequences. All further steps needed for the lower bound are proved in this manuscript.

The dependency chain is
\[
 \begin{gathered}
 \text{recurrent-ideal representatives and token upper bound}\\
 \Downarrow\\
 \text{cofinal synchronization and the two guards}\\
 \Downarrow\\
 \text{finite Cartesian-product embeddings}\\
 \Downarrow\\
 \text{ordinary ordinal-power amplification}\\
 \Downarrow\\
 \text{support-family classification and the finite-poset formula.}
 \end{gathered}
\]

In particular, the optional natural-product calculation for fixed lengths is not required to establish \cref{thm:main}. A reviewer can omit \cref{sec:exact} when checking the resolution of the original two-letter question.

\subsection{The critical points to check independently}

The proof has four particularly important boundaries. First, all embeddings are strict on ordinal positions; otherwise cofinal consumption fails. Second, a new support is not contained in any old support; otherwise a marker can be matched to an old block. Third, explicit periodic segments may begin inside canonical blocks, so exclusion must apply to the pure periodic segment rather than to its finite prefix. Fourth, full support requires an old proper support to build a positive limit-length guard.

Each boundary has a separate role in the proof. The negative controls show that removing the second or fourth safeguard, or omitting guards entirely, gives false order-reflection claims. The unary alphabet illustrates the resulting exceptional maximal order type.

\subsection{What has not been resolved here}

The argument does not compute $\ot(s_{\omega^3}(X))$, does not treat arbitrary infinite alphabets, and does not give general formulas at all countable length cutoffs. Its cofinal-block lemma exploits the fact that a target below $\omega^2$ has only finitely many $\omega$-blocks. At a larger cutoff, an image of order type $\omega$ can instead pass through infinitely many such blocks, so the same finite synchronization argument cannot simply be reused unchanged.

The formulas also do not classify the full isomorphism type of the transfinite-word order from $n$ and $|\J(X)|$. They identify its maximal order type. Two alphabets with the same ideal count can have different order-theoretic structure beyond that invariant.

Finally, the code is not a Lean formalization and is not presented as one. A formalization would need ordinal-indexed concatenation, the cofinal-block lemma, the two guards, and the classical maximal-order-type infrastructure. The finite token comparator and its tests provide a useful executable specification, but they are not a substitute for that development.

\section{Conclusion}

The finite generator upper bound for words below $\omega^2$ is sharp for every finite poset alphabet with at least two elements. The quantity that matters is the number of letters plus the number of nonempty recurrent ideals. Each proper recurrent ideal can be activated by a terminal-letter guard; the final full ideal can be activated only after an earlier proper ideal supplies a limit-length guard.

This gives the claimed exact value $\omega^{\omega^4}$ for the two-letter antichain, the value $\omega^{\omega^3}$ for the two-letter chain, and the general formula \eqref{eq:main}. It also explains why the universal-only family and the one-letter alphabet fall below the naive generator count. The solution proposed here is therefore not just a numerical lower bound: it identifies and isolates the structural mechanism responsible for sharpness.

\appendix
\section{A compact self-contained statement of the construction}

For quick reference, fix a nonempty finite poset $X$, an allowed support family $\F$, and a new nonempty ideal $D$ such that $D\nsubseteq E$ for all $E\in\F$. Set
\[
 g_D(u)=
 \begin{cases}
  ua,&D\neq X,\quad a\in X\setminus D,\\
  ua c_R,&D=X,\quad R\in\F\text{ proper},\quad a\in X\setminus R.
 \end{cases}
\]
The second line is defined only when the required $R$ exists. For every $m\in\N$, put
\[
 \Phi_m(\mathbf u)
  =g_D(u_0)c_D g_D(u_1)c_D\cdots c_D g_D(u_m).
\]
Then
\[
 \boxed{\quad
 \Phi_m(\mathbf u)\sub\Phi_m(\mathbf v)
 \quad\Longleftrightarrow\quad
 (\forall i\leq m)\ u_i\sub v_i.
 \quad}
\]
The proof is exactly synchronization plus nonleakage plus cancellation. If $\ot(\W_\F)=\HH_q$, this equivalence yields
\[
 \HH_{q+1}
 =\sup_{m<\omega}\HH_q^{m+1}
 \leq\ot(\W_{\F\cup\{D\}})
 \leq\HH_{q+1}.
\]
This is the complete lower-bound mechanism used in the article.

\clearpage
\section{Archive contents}

\begin{longtable}{@{}p{0.38\linewidth}p{0.56\linewidth}@{}}
\toprule
File & Purpose\\
\midrule
\endhead
\texttt{article.tex} & Complete editable source of the manuscript.\\
\texttt{article.pdf} & Compiled article.\\
\texttt{ordinal\_words.py} & Exact symbolic word comparator, normalization, ideals, guards, and tuple encodings.\\
\texttt{verify.py} & Exhaustive and seeded tests, including negative controls.\\
\texttt{verification\_results.json} & Recorded output of the default verification run.\\
\texttt{finite\_posets.csv} & The $50$ naturally labelled posets through size four, ideal counts, and predicted main-formula exponents.\\
\texttt{README.md} & Usage, scope, dependencies, and research-status notes.\\
\texttt{build.sh} & Three-pass PDF build command.\\
\bottomrule
\end{longtable}

\begin{thebibliography}{9}
\bibitem{Altman}
Harry Altman.
\emph{Bounding finite-image sequences of length $\omega^k$}.
\href{https://arxiv.org/abs/2409.03199v2}{arXiv:2409.03199v2},
10 March 2026; document dated 9 March 2026.
The explicit target is Example~3.15, printed page~10.

\bibitem{CP}
Alakh Dhruv Chopra and Fedor Pakhomov.
\emph{Well-quasi-orders on finite trees and transfinite sequences}.
\href{https://arxiv.org/abs/2602.09830v1}{arXiv:2602.09830v1},
10 February 2026.

\bibitem{DJP}
D.~H.~J. de Jongh and Rohit Parikh.
Well-partial orderings and hierarchies.
\emph{Indagationes Mathematicae} \textbf{39}(3) (1977), 195--207.
Also published in \emph{Proceedings of the Koninklijke Nederlandse Akademie van Wetenschappen, Series A} \textbf{80}(3).
\href{https://doi.org/10.1016/1385-7258(77)90067-1}{doi:10.1016/1385-7258(77)90067-1}.

\bibitem{Higman}
Graham Higman.
Ordering by divisibility in abstract algebras.
\emph{Proceedings of the London Mathematical Society}, series~3,
\textbf{2} (1952), 326--336.

\bibitem{Schmidt}
Diana Schmidt.
\emph{Well-partial orderings and their maximal order types}.
Habilitationsschrift, Heidelberg, 1979.
Reprinted in \emph{Well-Quasi Orders in Computation, Logic, Language and Reasoning},
Springer, 2020, pp.~351--391.
\end{thebibliography}
\end{document}
